\begin{lemma}[Scheide internal block-pair edges (3)]
Let $\gamma$ be a block color in a reduced Scheide configuration, put
$U_\gamma=\{u^1_\gamma,u^2_\gamma\}$ and
$Z_\gamma=V(T)\setminus U_\gamma$.  If
$\alpha\ne\gamma$ is missing at a vertex of $Z_\gamma$, then there is an
individual edge identity of color $\alpha$ whose endpoint pair is exactly
$\{u^1_\gamma,u^2_\gamma\}$.
\end{lemma}
\begin{proof}
Fix a block index $i$ and a color $\alpha\ne\gamma_i$, and abbreviate
\[
 X=\noderef{TreeSequenceVertexSet}(T),\qquad
 u_0=u^1_{\gamma_i},\quad u_1=u^2_{\gamma_i},\qquad
 Z=X\setminus\{u_0,u_1\}.
\]
Choose $z\in Z$ at which $\alpha$ is missing.  We keep edge identities
throughout the argument, even when two identities have the same endpoint
pair.

Unpack the extension in
\noderef{ScheideReducedConfiguration} to the underlying
\noderef{ScheideStructuralConfiguration}.  It supplies the integers
$p\geq2$ and $s$, the injectively indexed colors $\gamma_r$, the base
vertices $y_0,\ldots,y_{p-1}$, the base edges, and the two vertices and
two edges in each block.  Its position equalities say that these named
data, in precisely that order, are all vertices and edges of $T$.  Its
attachment clauses say that every block edge has endpoints
$\{y_q,u^j_{\gamma_r}\}$ for some base vertex $y_q$; its color clauses
give that edge color $\gamma_r$; its missing-color clauses give a base
vertex missing $\gamma_r$; and its base-avoidance clause says that every
colored base edge other than the first has color different from every
$\gamma_r$.  It also supplies a globally maximum triple.  Expanding
\noderef{GlobalMaximumTriple}, write its first component as the unchanged
$k$-triple $K$, its second as the Tashkinov certificate for $T$, and its
third as
\[
 |X|=t(G).                                                     \tag{1}
\]
By \noderef{ScheideGlobalMaximumOrder} applied to this globally maximum
triple, $X$ is elementary and closed in the senses of
\noderef{ColoredVertexSetElementary} and
\noderef{ColoredVertexSetClosed}.

The block-position equalities and vertex injectivity in
\noderef{TreeSequence} show that $u_0,u_1\in X$, that $u_0\ne u_1$, and
that neither equals any base vertex or a vertex in another block.  Since
$z\in Z$, it is distinct from both $u_0$ and $u_1$.  If $\alpha$ were
missing at either $u_j$, then the missing-color sets at $z$ and $u_j$
would meet in $\alpha$, contrary to elementariness.  Thus $\alpha$ is not
in $\noderef{MissingColors}(G,\varphi,u_j)$.  Since missing colors are,
by that definition, the complement of
\noderef{PresentColors}(G,\varphi,u_j)$, for each $j\in\{0,1\}$ there is
an individual edge identity $a_j$ such that
\[
 \noderef{Incident}(G,u_j,a_j),\qquad
 \varphi(a_j)=\alpha.                                         \tag{2}
\]

These incident $\alpha$-edges are unique.  Indeed, expansion of
\noderef{KTriple} gives the properness of $\varphi$ in the sense of
\noderef{IsProperEdgeColoring}.  Two distinct $\alpha$-colored identities
incident with the same $u_j$ would be adjacent by
\noderef{AdjacentEdges}, contradicting that properness.  By
\noderef{Incident}, the endpoint pair of $a_j$ contains $u_j$; write it
as
\[
 G.\operatorname{ends}(a_j)=\{u_j,w_j\}.
\]
The loopless clause in \noderef{LooplessMultigraph} gives $w_j\ne u_j$.
Moreover $w_j\in X$.  Otherwise $a_j$ would, by
\noderef{BoundaryEdges}, be a boundary edge of $X$.  But $\alpha$ is
missing at $z\in X$, so the closedness clause would forbid the
$\alpha$-colored boundary edge $a_j$, contrary to (2).

It remains to prove $w_0=u_1$.  Suppose for a contradiction that
$w_0\in Z$.  We first show $w_1\in Z$.  We already know $w_1\in X$ and
$w_1\ne u_1$.  If $w_1=u_0$, then $a_1$ is an $\alpha$-edge incident
with $u_0$.  It is not $a_0$: equality would make
$\{u_0,w_0\}=\{u_1,u_0\}$, whereas $w_0\in Z$ is distinct from $u_1$.
Thus $a_0$ and $a_1$ would be distinct adjacent edges of color $\alpha$,
contrary to properness.  Hence $w_1\ne u_0$, and so $w_1\in Z$.

We now construct the replacement ordered sequence $T'$.  Retain the
entire base, retain every block with index $r\ne i$ in increasing original
order, delete the two entries
\[
 (f^1_{\gamma_i},u_0),\qquad(f^2_{\gamma_i},u_1),
\]
and append
\[
 (a_0,u_0),\qquad(a_1,u_1).
\]
Thus, suppressing the unchanged base entries, the block part of $T'$ is
\[
 \mathop{\Vert}_{0\le r<s,\ r\ne i}
   (f^1_{\gamma_r},u^1_{\gamma_r},
    f^2_{\gamma_r},u^2_{\gamma_r}),
 \quad a_0,u_0,a_1,u_1,                                      \tag{3}
\]
where $\Vert$ means concatenation in increasing $r$.  We verify every
field of \noderef{TreeSequence} for (3).

First, (3) has $p+2(s-1)+2=p+2s$ vertices and one fewer edge.  The
length equality in the structural configuration therefore gives it the
same positive length as $T$.  Since $p\geq2$, at least one base edge is
retained; the first edge is the old first base edge and hence is still
$e_0$.  Every named vertex of $T$ occurs exactly once in (3): all vertices
except $u_0,u_1$ are retained once, and those two are appended once.
Consequently vertex distinctness follows from the old vertex-position
equalities and the old vertex injectivity.

For edge distinctness, the retained identities are pairwise distinct by
the edge injectivity of $T$.  Neither $a_j$ is a removed
$\gamma_i$-block edge, since (2), the block-color equation, and
$\alpha\ne\gamma_i$ give different colors.  Neither $a_j$ is the first
base edge: \noderef{KTriple} and
\noderef{IsOneEdgeUncolored} say that the latter is uncolored, whereas
(2) colors $a_j$.  A later base edge has, by its old
\noderef{TreeSequence} attachment and its structural position, both
ends among the base vertices; it cannot equal $a_j$, which is incident
with the nonbase vertex $u_j$.  A retained block edge has exact endpoints
$\{y_q,u^h_{\gamma_r}\}$ with $r\ne i$ and therefore cannot equal
$a_j$: old vertex injectivity makes $u_j$ distinct from both of those
endpoints.  Equivalently, if its color happened to be $\alpha$, the same
conclusion also follows by comparing these exact endpoint pairs; if it did
not, the color equations already separate the identities.  Finally
$a_0\ne a_1$, because their endpoint pairs contain respectively the
distinct removed vertices $u_0,u_1$, while their other ends $w_0,w_1$
both lie in $Z$.  This proves pairwise edge distinctness.

For the attachment field, every base edge retains its old attachment
witness because the base prefix is unchanged.  Each retained block edge
still has endpoints $\{y_q,u^h_{\gamma_r}\}$ by the structural
attachment clause, and every base vertex precedes every retained block
vertex in (3).  Finally (2) and the choices of $w_0,w_1$ give the exact
endpoint pairs $\{u_0,w_0\}$ and $\{u_1,w_1\}$.  Both $w_0,w_1$ lie in
$Z$, so both occur in the retained portion before the two appended
vertices.  These are valid earlier-vertex witnesses for the last two
edges.  Thus (3), with these data, is a tree-sequence $T'$.

Next verify \noderef{TashkinovTree} for $T'$.  There is no condition on
its first edge.  For each nonfirst base edge, use its original Tashkinov
witness; the whole base prefix and all earlier base vertices are
unchanged.  A retained block edge belonging to index $r\ne i$ has color
$\gamma_r$, and the structural missing-color clause supplies a base
vertex $y_q$ missing $\gamma_r$; that base vertex precedes the block edge
in $T'$.  Both appended edges have color $\alpha$ by (2), and the chosen
vertex $z\in Z$, which precedes both appended edges, misses $\alpha$.
These cases exhaust every nonfirst edge of $T'$, so $T'$ is Tashkinov.

The construction merely moves $u_0,u_1$ to the end and changes the two
edge identities paired with them.  Hence
\[
 \noderef{TreeSequenceVertexSet}(T')=X.                       \tag{4}
\]
The graph, distinguished edge, coloring, and $k$-triple $K$ are
unchanged.  Combining $K$, the preceding Tashkinov certificate, (4), and
(1), and then expanding \noderef{GlobalMaximumTriple}, proves that
$(T',e_0,\varphi)$ is globally maximum.

We next check directly that $\gamma_i$ is free for $T'$.  The structural
missing-color clause supplies a base vertex $y_q$ missing $\gamma_i$;
this vertex was retained, so
$\gamma_i\in\noderef{TreeMissingColors}(\varphi,T')$.  To see that the
color is unused, classify the tree edges in (3).  The first edge is
uncolored by the one-edge-uncolored clause of $K$ and in any event is
excluded from \noderef{TreeUsedColors}.  Every other base edge avoids
all block colors by the structural base-avoidance clause.  Every retained
block edge has color $\gamma_r$ for some $r\ne i$, and injectivity of the
indexed block colors gives $\gamma_r\ne\gamma_i$.  The last two edges
have color $\alpha\ne\gamma_i$.  Thus no colored nonfirst edge of $T'$
has color $\gamma_i$, so
$\gamma_i\notin\noderef{TreeUsedColors}(\varphi,T')$.  By the set
difference in \noderef{TreeFreeColors},
\[
 \gamma_i\in\noderef{TreeFreeColors}(\varphi,T').             \tag{5}
\]

It remains to transport the promised exit vertex, for which we do not
appeal to an unstated invariance lemma.  Expand
\noderef{DefectiveColors}: it is determined by $G$, $\varphi$, and the
vertex finset, so (4) makes the defective-color sets for $T$ and $T'$
equal.  Expand \noderef{ScheideDefectiveChainPath}: its only dependence
on the tree is through its vertex set, both in the condition that its
initial vertex lies in the set and in the condition that its terminal
vertex lies outside it; its simplicity and full-component clauses depend
only on $G$, $\varphi$, and the two colors.  Hence (4) makes every such
full-path predicate for $T$ equivalent to the corresponding predicate
for $T'$.  Now expand \noderef{ScheideExitVertex}.  Its tree-dependent
clauses are exactly membership in the tree vertex set, membership in the
just-identified defective-color set, the just-identified full-path
predicate, and the assertion that later vertex-nodes lie outside the same
tree vertex set.  Each clause is unchanged by (4).  Therefore a vertex is
a Scheide exit for $T$ if and only if it is one for $T'$.  Finally,
\noderef{ScheideExitVertices} filters the finite vertex universe by this
predicate, so
\[
 \noderef{ScheideExitVertices}(T,\varphi)
   =\noderef{ScheideExitVertices}(T',\varphi).                 \tag{6}
\]

The field ``every block color occurs at an exit'' in
\noderef{ScheideReducedConfiguration} now gives a vertex $v$ in the
left-hand side of (6) with
$\gamma_i\in\noderef{MissingColors}(G,\varphi,v)$.  By (6), $v$ is an
exit vertex of $T'$.  Apply \noderef{ScheideExitNoFree} to the globally
maximum triple just constructed.  It says that the missing-color set at
$v$ is disjoint from the free-color set of $T'$, contradicting the
membership of $\gamma_i$ in the former and (5) in the latter.

Thus $w_0\notin Z$.  We already proved $w_0\in X$ and looplessness gives
$w_0\ne u_0$; since $X$ is the disjoint union of $Z$ and
$\{u_0,u_1\}$, it follows that $w_0=u_1$.  The edge identity $a_0$ from
(2) therefore satisfies
\[
 \varphi(a_0)=\alpha,\qquad
 G.\operatorname{ends}(a_0)=\{u_0,u_1\},
\]
which is the required individual edge identity and endpoint equation.
\end{proof}
